\section{Related Work}

Our work connects three research threads: multi-dimensional trustworthiness evaluation, calibration and uncertainty quantification, and latent factor analysis in machine learning.

\subsection{Multi-Dimensional Trustworthiness Evaluation}

Recent frameworks conceptualize LLM trustworthiness as a multi-faceted construct. \citet{zhou2024trustrag} propose the Trust-RAG Compass with six dimensions: factuality, robustness, fairness, transparency, accountability, and privacy. This framework provides conceptual clarity but does not investigate cross-dimensional relationships. Similarly, TrustLLM \citep{sun2024trustllm} evaluates models across multiple trust dimensions without analyzing whether performance on one predicts performance on others.

Benchmark suites like HELM \citep{liang2023helm} and the Open LLM Leaderboard \citep{beeching2023openllm} report multi-benchmark scores, implicitly treating dimensions as independent. Our analysis of 4,561 models reveals that this independence assumption may be violated: cross-benchmark correlations of $\rho = 0.80$--$0.87$ suggest substantial shared variance.

\subsection{Calibration and Representation Stability}

\citet{khanmohammadi2025ccps} introduce CCPS (Calibrating via Perturbed Stability), achieving 55\% ECE reduction by leveraging internal representation stability. Their key insight---that perturbation-stable representations yield better calibration---aligns with our mechanism hypothesis. However, CCPS studies single-dimension calibration without examining cross-benchmark effects.

Our work extends this thread by hypothesizing that representation stability underlies not just calibration but the entire cross-benchmark correlation structure.

\subsection{Latent Factor Analysis in ML Evaluation}

In psychometrics, Spearman's $g$-factor explains positive correlations across cognitive tests. \citet{schumacher2024factor} analyze benchmark correlation structure, finding that model scale explains substantial shared variance. Our work extends their analysis by residualizing on scale and training recency before extracting latent factors, revealing structure beyond what scale alone explains.
